%! AP Statistics — Unit 3, Lesson 3.4 — Lesson Plan
\documentclass[10pt]{article}
\usepackage{apstats-article}
\usepackage{apstats-boxes}

\newcommand{\UnitNumberName}{Unit 3: Collecting Data \quad}
\newcommand{\LessonNumberName}{Lesson 3.4: Potential Problems with Sampling}

\begin{document}
\begin{center}
    {\Large\bfseries \CourseName: \SchoolYear\ (\MeetingLength) \\
    {\small\bfseries \UnitNumberName \LessonNumberName}}
\end{center}

\begin{tcolorbox}[colback=sky, colframe=navy, boxrule=0.9pt, arc=2mm,
    left=2mm, right=2mm, top=1.5mm, bottom=1.5mm]
    \textbf{Primary Objective:} Students will identify and explain the major sources
    of bias in sampling --- voluntary response, undercoverage, nonresponse, and
    response bias --- and recognize that non-random methods introduce these problems
    (Big Idea: DAT).
\end{tcolorbox}

\renewcommand{\arraystretch}{1.6}
\begin{skillbox}[Priority Ideas \& Skills]{goldbox}
\begin{minipage}[t]{0.48\linewidth}
    \textbf{AP Skill 1 --- Selecting Statistical Methods}
    \begin{itemize}
        \item Identify potential sources of bias in a sampling method (DAT-2.E,
              Skill 1.C).
        \item Explain how each bias type causes systematic over- or under-estimation.
        \item Distinguish the four main bias types by cause and mechanism.
    \end{itemize}
\end{minipage}
\hfill
\begin{minipage}[t]{0.48\linewidth}
    \textbf{Key Understandings}
    \begin{itemize}
        \item Bias: certain responses are systematically favored over others.
        \item Voluntary response bias: only those with strong opinions respond.
        \item Undercoverage: parts of the population are excluded from the sampling
              frame.
        \item Nonresponse: selected individuals decline to participate.
        \item Response bias: question wording, interviewer effects, or self-reporting
              distort answers.
    \end{itemize}
\end{minipage}
\end{skillbox}

\begin{skillbox}[{Vocabulary, Concepts, \& Theorems}]{greenbox}
\begin{minipage}[t]{0.48\linewidth}
    \begin{itemize}
        \item \textbf{Bias} --- a systematic tendency for the sample to misrepresent
              the population
        \item \textbf{Voluntary response bias} --- only volunteers (usually those with
              strong opinions) respond
        \item \textbf{Undercoverage bias} --- part of the population has a reduced
              chance of being included
    \end{itemize}
\end{minipage}
\hfill
\begin{minipage}[t]{0.48\linewidth}
    \begin{itemize}
        \item \textbf{Nonresponse bias} --- selected individuals who do not respond
              may differ systematically from those who do
        \item \textbf{Response bias} --- survey instrument or context causes
              systematically inaccurate answers (question wording, leading questions,
              self-reporting)
        \item \textbf{Convenience sample} --- individuals chosen because they are
              easy to reach; introduces potential for all bias types
    \end{itemize}
\end{minipage}
\end{skillbox}

\begin{skillbox}[Activate Prior Knowledge \& Spiral Review]{sky}
\begin{minipage}[t]{0.48\linewidth}
    \textbf{Prior Knowledge Check (3--4 min)}
    \begin{itemize}
        \item Recall: ``From Lesson 3.3, what makes a simple random sample
              trustworthy?'' (every group has equal chance)
        \item Transition: ``But even well-intentioned surveys can go wrong. Today we
              look at the ways sampling can fail.''
    \end{itemize}
\end{minipage}
\hfill
\begin{minipage}[t]{0.48\linewidth}
    \textbf{Connections Forward}
    \begin{itemize}
        \item Bias in experiments (confounding) $\to$ Lesson 3.5
        \item How experiments eliminate bias with random assignment $\to$ Lessons 3.5--3.6
        \item Interpreting survey results on the AP exam $\to$ Unit 9
    \end{itemize}
\end{minipage}
\end{skillbox}

\begin{skillbox}[Hook \& Lesson]{sky}
\begin{minipage}[t]{0.48\linewidth}
    \textbf{Hook (5--7 min)}\\
    Show this survey question: \textit{``Don't you agree that the new cafeteria food
    is terrible?''}\\[4pt]
    Ask: ``What's wrong with this question?'' Then show a second version: ``How would
    you rate the new cafeteria food? Excellent / Good / Fair / Poor.''\\[4pt]
    Transition: \textit{``Even if we use a perfect random sample, the question itself
    can bias the results.''}
\end{minipage}
\hfill
\begin{minipage}[t]{0.48\linewidth}
    \textbf{Lesson Flow (15--18 min)}
    \begin{enumerate}
        \item Define bias: systematic error, not random.
        \item Voluntary response: call-in polls, online votes --- who responds?
        \item Undercoverage: landline phone surveys miss cell-phone-only households.
        \item Nonresponse: people who are selected but refuse to participate may have
              systematically different opinions.
        \item Response/question-wording bias: leading questions, sensitive topics,
              self-reporting errors.
        \item Non-random sampling methods (convenience, judgment) introduce all of
              these problems without a chance to correct them.
    \end{enumerate}
\end{minipage}
\end{skillbox}

\begin{skillbox}[Explicit Instruction \& Active Monitoring]{sky}
\begin{minipage}[t]{0.48\linewidth}
    \textbf{Direct Instruction (10 min)}
    \begin{itemize}
        \item Walk through 4 flawed surveys. For each: (a) identify the bias type,
              (b) explain the mechanism (which direction does it push the estimate?),
              (c) suggest a fix.
        \item Emphasize: ``bias'' on the AP exam means \emph{systematic}
              over- or under-estimation --- not just ``unfair.''
    \end{itemize}
\end{minipage}
\hfill
\begin{minipage}[t]{0.48\linewidth}
    \textbf{Active Monitoring}
    \begin{itemize}
        \item Listen for: students identifying bias without explaining the direction
              of distortion.
        \item Cold-call: ``If only people who hate the policy call in, is the estimate
              too high or too low? Why?''
        \item Check activity work for precision: ``nonresponse'' is not the same as
              ``no one responded.''
    \end{itemize}
\end{minipage}
\end{skillbox}

\begin{skillbox}[Group Work \& Differentiation]{redbox}
\begin{minipage}[t]{0.48\linewidth}
    \textbf{Partner Activity (8--10 min)}\\
    Four flawed survey descriptions. Partners:
    \begin{enumerate}
        \item Identify the type of bias (may be more than one).
        \item Explain how the bias would distort the estimate (too high / too low /
              in a specific direction).
        \item Suggest a redesign that would reduce the bias.
    \end{enumerate}
\end{minipage}
\hfill
\begin{minipage}[t]{0.48\linewidth}
    \textbf{Differentiation}
    \begin{itemize}
        \item \textit{Support}: Bias identification card --- definition + one example
              for each type.
        \item \textit{Extension}: Find a real survey in the news. Identify at least one
              potential source of bias and explain how it could affect the results.
        \item \textit{ELL}: Sentence frame --- ``This survey has \underline{\hspace{1cm}}
              bias because \underline{\hspace{2cm}}, which means the estimate is probably
              too \underline{\hspace{1cm}}.''
    \end{itemize}
\end{minipage}
\end{skillbox}

\begin{skillbox}[Individual Work \& Assessment]{redbox}
\begin{minipage}[t]{0.48\linewidth}
    \textbf{Exit Ticket (5 min)}\\
    Scenario: \textit{``A researcher mails surveys to 1{,}000 randomly selected
    adults to ask about their annual income. Only 120 return the survey, and most
    are from high-income households.''}
    \begin{enumerate}
        \item Identify the type(s) of bias present.
        \item Explain how this bias distorts the estimate of average income.
        \item Was the original selection (1{,}000 random adults) biased?
    \end{enumerate}
\end{minipage}
\hfill
\begin{minipage}[t]{0.48\linewidth}
    \textbf{AP-Style Check}\\
    \textit{``A political poll asks: `Are you in favor of wasteful government
    spending?' Which type of bias is most likely introduced by this question?''}
    \begin{enumerate}[label=(\Alph*)]
        \item Undercoverage bias
        \item Voluntary response bias
        \item Nonresponse bias
        \item Response bias (question wording)
    \end{enumerate}
    Answer: (D)
\end{minipage}
\end{skillbox}

\begin{skillbox}[Reinforcement \& Extension]{goldbox}
\begin{minipage}[t]{0.48\linewidth}
    \textbf{Homework / Practice}
    \begin{itemize}
        \item For 4 flawed surveys: identify the bias type, explain the direction of
              distortion, and suggest an improvement.
        \item Reflection: \textit{``Can you have a large, biased sample and a small,
              unbiased sample at the same time? Which is better?''}
    \end{itemize}
\end{minipage}
\hfill
\begin{minipage}[t]{0.48\linewidth}
    \textbf{Extension / Preview}
    \begin{itemize}
        \item \textit{Extension}: Write two versions of the same survey question ---
              one with response bias and one without. Explain how the wording
              affects responses.
        \item \textit{Preview}: Lesson 3.5 shifts to experiments. Think about how
              random \emph{assignment} (not selection) eliminates confounding.
    \end{itemize}
\end{minipage}
\end{skillbox}

\end{document}
